%% LyX 2.3.3 created this file.  For more info, see http://www.lyx.org/.
%% Do not edit unless you really know what you are doing.
\documentclass[english]{article}
\renewcommand{\sfdefault}{cmss}
\usepackage[T1]{fontenc}
\usepackage[latin9]{inputenc}
\usepackage{amsmath}
\usepackage{graphicx}
\usepackage{babel}
\begin{document}

\section*{Transport equation}

The convection-diffusion equation for the transport of temperature
\textit{T} is

\begin{equation}
\frac{\partial T}{\partial t}=-U\frac{\partial T}{\partial x}+k\frac{\partial}{\partial x}(\frac{\partial T}{\partial x})+S
\end{equation}
where \textit{U} is velocity and \textit{S} a source term. For a non
existing convection case, Equation (1) becomes the diffusion Equation

\begin{equation}
\frac{\partial T}{\partial t}=k\frac{\partial}{\partial x}(\frac{\partial T}{\partial x})+S
\end{equation}
that, for a steady-state, is 

\begin{equation}
0=k\frac{\partial}{\partial x}(\frac{\partial T}{\partial x})+S
\end{equation}
These equations are solved by a Finite Volume Method (FVM) and by
a Finite Difference Method (FDM). 

\section*{Finite Volume Method }

In the Finite Volume Method, values for the above diferential equations
are calculated at discrete places on a grid of volumes, shown in Figure
1,
\begin{center}
\includegraphics{figure_1}
\par\end{center}

\begin{center}
\textbf{Figure 1}: Temperature at center of cell, and on left and
right cell borders.
\par\end{center}

\noindent where temperature at the center of a cell is \textit{T}\textit{\scriptsize{}P},
\textit{T}\textit{\scriptsize{}r} e \textit{T}\textit{\scriptsize{}l}
is temperature at the border between cells, right and left, \textit{T}\textit{\scriptsize{}R}
e \textit{T}\textit{\scriptsize{}L} is temperature at the right and
left cell centers. Integrating Equation 3 over a cell volume, and
considering the rate of accumulation over the volume \textit{V} equal
to the flow across the surfaces of the control volume, Equation (6),
where \textit{n} is the unit normal vector pointing out of the control
volume and \textit{A} is the cross sectional area of the volume.

\begin{equation}
\int\left[k\frac{\partial}{\partial x}(\frac{\partial T}{\partial x})+S\right]dV=0
\end{equation}

\begin{equation}
\int\left[\frac{\partial}{\partial x}(k\frac{\partial T}{\partial x})dV\right]+SV=0
\end{equation}

\begin{equation}
\int(k\frac{\partial T}{\partial x}n)dA+SV=0
\end{equation}

\noindent Equation (6) is written in terms of the flux leaving the
right face \textit{r} minus the flux entering the left face \textit{l}
in Figure (1)

\begin{equation}
(kA\frac{\partial T}{\partial x})_{r}-(kA\frac{\partial T}{\partial x})_{l}+SV=0
\end{equation}

\noindent Equation (7) for the temperature gradient at the boundaries
is dicretized in terms of the temperature at the cell centers, where
$\Delta x$ is the distance between cell centers

\begin{equation}
kA\frac{T_{R}-T_{P}}{\Delta x}-kA\frac{T_{P}-T_{L}}{\Delta x}+SV=0
\end{equation}

\noindent An individual set of equations is made for each cell.

\section*{Finite Difference Method}

The diffusion equation (2) is approximated by finite differences with 

\begin{equation}
\frac{T_{i}^{n+1}-T_{i}^{n}}{\varDelta x}=k\frac{T_{i-1}^{n}-2T_{i}^{n}+T_{i+1}^{n}}{\varDelta x^{2}}+S_{i}^{n}
\end{equation}


\section*{References}
\begin{enumerate}
\item Computational Fluid Dynamics Fundamentals Course. A. Wimshurst. 2019. 
\item An Introduction to Computational Fluid Dynamics: The Finite Volume
Method. H. Versteeg, W. Malalasekera. 2007.
\item Finite Difference Computing with PDEs. A Modern Software Approach.
H. Langtangen, S. Linge. 2016.
\end{enumerate}

\end{document}
